% Section 8 -- Limitations (shortened to the two main limitations, 2026-11-24, user decision).
\section{Limitations}\label{sec:limitations}

The main limitation is that the class-pure sites were constructed from cohorts in which the sites hold both classes. Drawing both roles at random from the same cohort controls part, but not all, of the difference between case sites and control sites, and does not guarantee (A1$'$), because a site can take a role only if it has enough participants of that class. In real collections, cases from a clinic and controls from a school may also differ in age, medication or motion, a bias that none of our methods removes; the natural collections show that the evaluation bias is real, but not whether bridging removes it there. The constructed designs were also small, so LSPO remained imprecise (mean absolute error $0.09$ to $0.19$ across the two ABIDE cohorts).

The theory assumes additive site effects (A0) and linear scores. Bridging and within-site training remove a site shift common to both classes but cannot address a site-by-diagnosis interaction, such as a scanner that distorts patients' signals differently. Proposition~\ref{prop:leak} was formulated after the experiments, and our L3 implementation is shift-invariant only approximately.
